% =====================================================================
% El dato crudo: cinco instituciones, veintisiete fuentes
% Parte I. El dato
%
% ENCARGO: que se midio, donde y con que. 20 medidores mas 7 inversores.
% Cobertura temporal y justificacion del horizonte de 6.144 h.
%
% Reglas del documento:
% - Una idea = una subseccion = una figura.
% - Toda figura declara su procedencia con.
% - Toda cifra sale del canon de agosto; ver GUIA.md.
% - M1 y M3 se reportan siempre en paralelo.
% - Ningun superlativo sin releer el CSV de la frontera correcta.
% =====================================================================
\section{Las fuentes de medición: cinco instituciones, veintisiete equipos}
\label{sec:dato}
\justifying

Todo lo que este documento afirma descansa sobre una base de medidas
tomadas en cinco instituciones de San Juan de Pasto entre abril y
diciembre de 2025. Antes de describir qué se hizo con esas medidas hay que precisar
qué son, cuántas hay y de qué calidad, porque la solidez de
cualquier conclusión posterior está delimitada por la de este punto de
partida.

\subsection{Las cinco instituciones}
\label{sub:dato-instituciones}

La comunidad estudiada la forman cinco entidades del municipio de Pasto,
todas ellas con sistema fotovoltaico instalado y medición avanzada: la
Universidad de Nariño, la Universidad Mariana, la Universidad Cooperativa
de Colombia (campus Pasto), el Hospital Universitario Departamental de
Nariño E.S.E. y la Universidad CESMAG. A lo largo del documento se
las nombra de forma abreviada como Udenar, Mariana, UCC, HUDN y CESMAG, y
ese orden se conserva en todas las tablas y figuras.

Estas cinco son las entidades que el proyecto MTE instrumentó y de las
que, en consecuencia, existen series numéricas de generación y de consumo.
La comparación se construye alrededor de ellas.

\begin{guardado}
\begin{trampabox}
Que el conjunto venga dado tiene una consecuencia que reaparece en el
Capítulo~\ref{sec:esc} y que no conviene descubrir allí: son cinco, y
cinco es un número pequeño con efectos normativos. El artículo 20 de la
Resolución CREG 101 072 de 2025 reserva su tratamiento más favorable, el
del numeral 1, a las comunidades que cumplen tres condiciones, y la
tercera exige que el Porcentaje de Distribución de Excedentes sea
«inferior al 10\% para cada uno de los usuarios». Ese porcentaje lo
acuerdan los propios integrantes, pero el artículo 19 obliga a que sume
exactamente el 100\%, de modo que la elección no es libre: si todos los
usuarios quedaran por debajo del 10\%, la suma sería inferior a 10 veces
el número de usuarios, y alcanzar el 100\% exigiría entonces 11
participantes o más. Con cinco, la condición resulta aritméticamente
inalcanzable, y no por una característica de esta comunidad sino por su
tamaño. Tampoco hubo margen para ampliar el conjunto ni para controlar su
composición, de modo que ninguna conclusión de este trabajo se extiende
sin más a comunidades de otro tamaño o de otra mezcla de usos.
\end{trampabox}
\end{guardado}

El sistema instalado en la Universidad de Nariño ocupa el Bloque
Sur de la sede Torobajo. El de la Universidad Mariana está en el Campus Deportivo Alvernia, en el sector
occidental de la ciudad, y de ahí que los equipos que lo miden lleven ese
nombre. Los de la Universidad Cooperativa de Colombia y el HUDN sí están
en sus sedes principales, en el sector Torobajo y en el Parque Bolívar
respectivamente. El de la Universidad CESMAG es el único fuera del casco
urbano: está en el Campus Universitario San Damián, en el
corregimiento de Catambuco. La Figura~\ref{fig:dato-mapa} sitúa
las cinco sobre el territorio.

\begin{figure}[H]
    \centering
    \includegraphics[width=0.98\textwidth]{figuras/f2_01_mapa_ubicacion.pdf}
    \caption[Localización de las cinco instituciones]{Localización de las
    instituciones aliadas del proyecto MTE en el municipio de Pasto. El
    área sombreada en rosa es el perímetro urbano; el fondo verde, el
    modelo digital de elevación.
    }
    \label{fig:dato-mapa}
\end{figure}

De ese conjunto interesa una característica, que no se supone sino que se
mide: las cinco entidades no comparten rutina de consumo. Cuatro son
instituciones educativas con jornada diurna y calendario académico,
mientras que la quinta es un hospital que no cierra nunca. Esa disparidad
es la que hace que en una misma hora unas tengan excedente y otras
déficit, que es la premisa material de cualquier intercambio. El
Capítulo~\ref{sec:prep} la muestra medida, institución por institución.

\begin{guardado}
\begin{detallebox}
Son dos proximidades distintas y se confunden con facilidad. La
Figura~\ref{fig:dato-mapa} acredita la \emph{geográfica}, y es
considerable: cuatro de las cinco caben en unos pocos kilómetros. De ahí
no se sigue la \emph{eléctrica}, que es la que importaría para hablar de
un mismo circuito o de una misma subestación, y que el proyecto no
documenta. Este trabajo no la afirma. Tampoco haría falta: el modelo del
Capítulo~\ref{sec:modelo} no representa la red de distribución ni calcula
flujos de potencia, de modo que la distancia eléctrica entre participantes
no interviene en ningún resultado.

Tampoco comparten comercializador. La mayoría de las cinco tiene contrato
con un comercializador distinto del operador de red, y aun así este
trabajo liquida las cinco contra la tarifa publicada por CEDENAR. Es una
decisión tomada por disponibilidad de la fuente, y el
Capítulo~\ref{sec:precios} la declara junto con su alcance.
\end{detallebox}
\end{guardado}

Sí interviene, en cambio, una asimetría regulatoria que el
Capítulo~\ref{sec:precios} desarrolla: dos de las cinco, Udenar y el HUDN,
son entidades públicas, y las otras tres, privadas. Las dos categorías
pagan tarifas distintas por la misma energía.

\subsection{Veintisiete fuentes, dos magnitudes}
\label{sub:dato-fuentes}

Cada institución está instrumentada con cuatro medidores eléctricos y al
menos un inversor fotovoltaico, y todos registran con resolución
de dos minutos. El inventario de los equipos arroja
27 dispositivos: 20 medidores y 7 inversores. Estos últimos
están repartidos de forma desigual, porque Udenar tiene tres y las demás
uno cada una. Hay además cinco estaciones meteorológicas instaladas, una por
institución, que este trabajo no emplea.

De cada equipo se toma una sola magnitud, y conviene subrayar cuánto se
descarta al hacerlo. Un medidor registra 54
magnitudes por lectura, entre ellas la corriente y la tensión de cada fase,
la frecuencia, el factor de potencia y la distorsión armónica; de todas
ellas el modelo usa una, la potencia activa total (kW). Un inversor
registra 18, entre ellas la potencia y la tensión del lado de
corriente continua; el modelo usa una, la potencia de corriente alterna
entregada (W).

La elección merece sustento, porque el mismo medidor ofrece otras dos
potencias, la reactiva y la aparente. Lo que el mercado intercambia es
energía, y la energía es la integral de la potencia en el tiempo,

\begin{equation}
 E \;=\; \int_{t}^{t+\Delta t} p(\tau)\,\mathrm{d}\tau
   \;\simeq\; \overline{P}\,\Delta t,
 \label{eq:dato-energia}
\end{equation}

de modo que la magnitud buscada es aquella cuya integral da los kilovatios
hora que después se facturan. De las tres potencias que el equipo mide,
solo la activa lo cumple, porque es la única que transfiere energía neta a
la carga. La reactiva oscila entre la fuente y la carga, y su
transferencia a lo largo de un ciclo se anula, de modo que su integral
sobre la hora no es energía consumida. La aparente combina las dos en una
sola magnitud y por eso tampoco sirve: no permite separar la parte que sí
se consume. Quedarse con la activa no sale gratis, porque la reactiva
también se factura, y el Anexo~\ref{sec:no-representado} mide ese cargo y
muestra que no altera el ordenamiento.

Por la misma razón se toma del inversor la potencia del lado de corriente
alterna y no la del lado de corriente continua, que el equipo también
registra. Lo que la institución puede consumir o vender es lo que sale del
inversor ya convertido, después de las pérdidas de conversión, mientras
que la lectura del lado continuo describe lo que entrega el arreglo
fotovoltaico y sobrestimaría la generación disponible.

Las Tablas~\ref{tab:dato-var-medidor} y~\ref{tab:dato-var-inversor} recogen todas las variables que entrega cada clase de equipo, agrupado por
familia y con la unidad de cada magnitud.

\begingroup
\footnotesize
\begin{longtable}{@{}>{\raggedright\arraybackslash}p{9.6cm} l@{}}
\caption{Las 54 magnitudes que entrega un medidor, por
familia. La marca de tiempo no se cuenta, porque no es una magnitud
medida.}\label{tab:dato-var-medidor}\\
\toprule
\textbf{Variable} & \textbf{Unidad} \\
\midrule
\endfirsthead
\multicolumn{2}{@{}l}{\itshape Tabla \thetable, continuación}\\
\toprule
\textbf{Variable} & \textbf{Unidad} \\
\midrule
\endhead
\midrule
\multicolumn{2}{r@{}}{\itshape continúa en la página siguiente}\\
\endfoot
\bottomrule
\endlastfoot
    \multicolumn{2}{@{}l}{\textit{Estado de la red}} \\*
    Frecuencia                                        & Hz \\
    Estado de salida & --- \\
    \addlinespace
    \multicolumn{2}{@{}l}{\textit{Corriente}} \\*
    Corriente de fase (A, B, C)                       & A \\
    Corriente máxima de fase (A, B, C)                & A \\
    Corriente de neutro                               & A \\
    \addlinespace
    \multicolumn{2}{@{}l}{\textit{Tensión}} \\*
    Tensión de fase (A, B, C)                         & V \\
    \addlinespace
    \multicolumn{2}{@{}l}{\textit{Potencia activa}} \\*
    \textbf{Potencia activa total} \hfill \textbf{(la que usa el modelo)}
                                                      & kW \\
    Potencia activa de fase (A, B, C)                 & kW \\
    Demanda de potencia activa vigente                & kW \\
    Demanda máxima de potencia activa                 & kW \\
    Potencia activa acumulada                         & kW \\
    \addlinespace
    \multicolumn{2}{@{}l}{\textit{Potencia reactiva}} \\*
    Potencia reactiva total                           & kvar \\
    Potencia reactiva de fase (A, B, C)               & kvar \\
    \addlinespace
    \multicolumn{2}{@{}l}{\textit{Potencia aparente}} \\*
    Potencia aparente total                           & kVA \\
    Potencia aparente de fase (A, B, C)               & kVA \\
    Demanda de potencia aparente vigente              & kVA \\
    Demanda máxima de potencia aparente               & kVA \\
    Potencia aparente acumulada                       & kVA \\
    Registro de potencia aparente, palabras baja y alta & --- \\
    \addlinespace
    \multicolumn{2}{@{}l}{\textit{Factor de potencia}} \\*
    Factor de potencia total                          & --- \\
    Factor de potencia de fase (A, B, C)              & --- \\
    Factor de potencia en la máxima aparente importada & --- \\
    \addlinespace
    \multicolumn{2}{@{}l}{\textit{Registros de energía acumulada}} \\*
    Energía activa importada, palabras baja y alta    & kWh \\
    Energía activa exportada, palabras baja y alta    & kWh \\
    Energía reactiva positiva, palabras baja y alta   & kvarh \\
    Energía reactiva negativa, palabras baja y alta   & kvarh \\
    \addlinespace
    \multicolumn{2}{@{}l}{\textit{Calidad de la onda}} \\*
    Distorsión armónica de corriente (A, B, C)        & \% \\
    Distorsión armónica de tensión (A, B, C)          & \% \\
    Distorsión de demanda de corriente (A, B, C)      & \% \\
\end{longtable}
\endgroup

\begin{table}[H]
    \centering
    \footnotesize
    \caption{Las 18 variables que entrega un inversor, por familia.}
    \label{tab:dato-var-inversor}
    \begin{tabular}{@{}>{\raggedright\arraybackslash}p{9.6cm} l@{}}
    \toprule
    \textbf{Variable} & \textbf{Unidad} \\
    \midrule
    \multicolumn{2}{@{}l}{\textit{Lado de corriente continua}} \\*
    Potencia                                          & W \\
    Corriente                                         & A \\
    Tensión                                           & V \\
    \addlinespace
    \multicolumn{2}{@{}l}{\textit{Lado de corriente alterna, potencia}} \\*
    \textbf{Potencia activa entregada} \hfill \textbf{(la que usa el modelo)}
                                                      & W \\
    Potencia aparente                                 & VA \\
    Potencia reactiva                                 & var \\
    Factor de potencia         & \% \\
    \addlinespace
    \multicolumn{2}{@{}l}{\textit{Lado de corriente alterna, medidas eléctricas}} \\*
    Frecuencia                                        & Hz \\
    Corriente total                                   & A \\
    Corriente de fase (A, B, C)                       & A \\
    Tensión de fase a neutro (A, B, C)                & V \\
    Tensión entre fases (AB, BC, CA)                  & V \\
    \bottomrule
    \end{tabular}
\end{table}

\begin{guardado}
\begin{detallebox}
Que el modelo consuma 2 de las 72 magnitudes disponibles no
es un desperdicio sino una consecuencia de su formulación: el mercado que
se simula intercambia energía, y para eso bastan la que cada institución
consume y la que produce. Las variables restantes describen la
\emph{calidad} de esa energía, que este trabajo no modela. Quedan
disponibles para extensiones que sí la consideren, y las estaciones
meteorológicas, para trabajos que quieran predecir la generación en lugar
de tomarla medida.
\end{detallebox}
\end{guardado}

\subsubsection*{Las magnitudes de cada equipo}

Hasta aquí los equipos se han citado con el nombre que les da la plataforma,
esto es, Medidor 1 o Medidor 3, y ese nombre no dice a qué circuito se
conectaron. El inventario de instalación del proyecto sí lo dice: cruza
cada equipo con el circuito que mide, el tablero donde está montado, el
transformador de corriente y la escala de conversión. La
Tabla~\ref{tab:dato-inventario} recoge de ahí los dos medidores que el
modelo emplea en cada institución. De ese cruce sale el nombre de las dos
fronteras de medición: las nombra el circuito que miden, no el número del
equipo.

\begin{table}[H]
    \centering
    \small
    \caption{Circuito medido por los dos equipos que el modelo emplea en
    cada institución.}
    \label{tab:dato-inventario}
    \begin{tabular}{@{}l >{\raggedright\arraybackslash}p{4.3cm}
                        >{\raggedright\arraybackslash}p{4.3cm}@{}}
    \toprule
    \textbf{Institución} & \textbf{Medidor 1 (M1)} &
    \textbf{Medidor 3 (M3)} \\
    \midrule
    Udenar   & Circuito principal, tablero de la subestación del bloque Sur
             & Circuito secundario, piso 1A \\
    \addlinespace
    Mariana  & Circuito de inyección, totalizador principal del campus
               Alvernia
             & \emph{No existe.} Se aproxima con el Medidor 1 escalado \\
    \addlinespace
    UCC      & Circuito de inyección, totalizador principal del bloque A
             & Totalizador del piso 2 \\
    \addlinespace
    HUDN     & Circuito de inyección principal, subestación del piso 3
             & Alimentación ininterrumpida de ginecología \\
    \addlinespace
    CESMAG   & Circuito principal, totalizador del bloque B
             & Circuito principal, totalizador del bloque A \\
    \bottomrule
    \end{tabular}
\end{table}

De cada institución el modelo lee uno de esos dos medidores, y el trabajo
entero se recorre dos veces, una con cada elección. Los dos recorridos se
llaman en adelante M1 y M3, por el número del medidor que cada uno emplea,
y el Capítulo~\ref{sec:frontera} desarrolla en qué se diferencian y por
qué se reportan siempre los dos.

\begin{guardado}
\begin{detallebox}
No todas las fuentes instrumentadas se usan. De los 20 medidores, el
modelo lee un equipo por institución en cada frontera, de modo que 9
participan y 11 quedan como auxiliares. De los 7 inversores, 5
definen la generación que ve el modelo, uno por institución, y los otros 2
son de Udenar: intervienen en la reconstrucción de su demanda, y uno de
ellos sirve además de referencia para extender hacia atrás la serie del
inversor del proyecto, que no entró en servicio hasta septiembre. La razón de no sumar todo lo
disponible se explica en el Capítulo~\ref{sec:prep}: los equipos miden
puntos distintos del circuito, y sumarlos mezclaría totalizadores con
ramales.
\end{detallebox}

\subsection{Las magnitudes excluidas}
\label{sub:dato-reactiva}

Que la energía reactiva no entregue trabajo no significa que no se pague.
La regulación colombiana cobra el consumo reactivo inductivo que supera la
mitad de la energía activa entregada en el mismo periodo horario, y lo
factura como si fuera activa dentro de los cargos por uso de las redes.
Medido sobre el horizonte, ese exceso aparece en el \pct{15,0} de las
horas con consumo del circuito principal y en el \pct{65,7} de las del
secundario, y equivale a \uni{1865016}{COP} y \uni{1210145}{COP}. La
medición no procede de los contadores de energía reactiva del equipo, que
llegan en cero durante todo el horizonte, sino de integrar hora a hora la
potencia reactiva, del mismo modo que se hace con la activa. La
primera cifra supera el margen que separa a los dos mecanismos mejor
situados en esa frontera. Aun así el ordenamiento no cambia, porque es un
cargo que todos los mecanismos pagan sobre el mismo consumo, y el efecto
propio del mercado sobre él vale el \pct{6,8} de ese margen. El
Anexo~\ref{sec:no-representado} recoge la medición completa y las razones
por las que el trabajo se hace solo sobre energía activa.
\end{guardado}

\subsection{Cobertura temporal de las fuentes}
\label{sub:dato-gantt}

La segunda pregunta después de «qué se midió» es «desde cuándo». Los
equipos no entraron en servicio a la vez, y esa asincronía condiciona el
horizonte del estudio. La Figura~\ref{fig:dato-gantt} sitúa cada fuente
sobre el eje temporal y deja ver el escalonamiento con que la comunidad
se fue instrumentando.

\begin{figure}[H]
 \centering
 \includegraphics[width=0.95\textwidth]{figuras/f2_02_gantt_fuentes.png}
 \caption[Cobertura temporal de las 27 fuentes]{Periodo cubierto por
 cada una de las 27 fuentes instrumentadas, agrupadas por institución. La
 banda vertical es el horizonte del estudio y el anillo señala la fuente
 cuyo primer registro fija su inicio.}
 \label{fig:dato-gantt}
\end{figure}

\subsection{Delimitación del horizonte de estudio}
\label{sub:dato-horizonte}

El inicio del horizonte no se eligió: viene impuesto. Un mercado entre
cinco participantes no puede simularse mientras falte uno, de modo que la
primera fecha utilizable es aquella en que empieza a registrar la última de
las fuentes esenciales. Son esenciales 14 de las 27 fuentes instrumentadas:
el medidor de cada frontera en cada institución, nueve en total, porque el
de Mariana sirve a las dos, y la fuente que provee la generación de cada
una desde el primer día.

En cuatro instituciones esa fuente es su único inversor. En Udenar es el
inversor de referencia con el que se extiende la serie del inversor del
proyecto, que no registra hasta septiembre. En la
Figura~\ref{fig:dato-gantt} las esenciales son las barras en color, y el
anillo señala la última en registrar.

Esa última es el inversor del HUDN, que empieza a registrar el 4 de abril
de 2025 a las 05:58. Que sea un inversor y no un medidor no es un detalle
menor: dentro del HUDN los dos medidores empiezan la víspera, de modo que
lo que faltaba para completar la comunidad era la generación. Quedan fuera
de la cuenta los otros dos equipos de Udenar: el inversor del proyecto,
cuya serie se reconstruye hacia atrás y por tanto no condiciona el
arranque, y el segundo Fronius, que solo interviene en la reconstrucción
de la demanda. Contar el primero señalaría como causa del arranque a un
equipo que empieza en septiembre y que no lo causa.

De ahí esa fecha como primer día del horizonte, el primero en que las cinco
instituciones tienen registro en sus medidores y en su inversor. Se cierra
el 16 de diciembre del mismo año, lo que da \num{256} días y \num{6144}
horas.

\subsection{Completitud del registro por fuente}
\label{sub:dato-cobertura}

Queda la tercera pregunta: dentro de ese horizonte, ¿cuántas horas faltan?
La respuesta debe darse por separado para cada clase de equipo, porque
medirlas igual sería un error de método: contar la noche como dato faltante
en un inversor confundiría la ausencia de sol con la ausencia de medición.
La Figura~\ref{fig:dato-cobertura} separa las dos clases y cuenta, en cada
fuente, las horas sin registro. El denominador es el horizonte y no la
serie completa del equipo, que llega hasta 2026: se cuentan solo las horas
del horizonte en que la fuente ya estaba registrando.

\begin{figure}[H]
 \centering
 \includegraphics[width=0.95\textwidth]{figuras/f2_03_cobertura_fuentes.png}
 \caption[Horas sin registro por fuente]{Horas sin registro de cada
 fuente. Los medidores se evalúan sobre las 24 horas del día; los
 inversores, solo sobre la franja diurna, que va de las 06:00 a las 19:00.}
 \label{fig:dato-cobertura}
\end{figure}

Medidos contra las \num{6144} horas del horizonte, dos equipos de Udenar
parecerían defectuosos: el Medidor 4 marcaría el \pct{88.7} y el Inversor
MTE, el \pct{39.3} de las horas de sol. Ninguno de los dos falla. El
primero se instaló el 29 de abril y el segundo entró en septiembre, de modo
que a los dos se les contaban como huecos horas en que el equipo todavía no
existía.

Descontadas esas horas, las 27 fuentes caen entre el \pct{96.7} y el
\pct{98.8} y no queda excepción que acotar. El peor caso son las \num{204}
horas que le faltan al Medidor 4 de Cesmag, que es auxiliar; de los equipos
que el modelo emplea, a ninguno le faltan más de \num{126}. Son pocas, y
son exactamente las que el capítulo siguiente tiene que rellenar.

\subsection{Alcance y límites del inventario}
\label{sub:dato-alcance}

Antes de pasar al tratamiento es importante delimitar el alcance de esta base de
datos, porque determina lo que el resto del documento puede afirmar.

Se tienen nueve meses continuos de medición horaria sobre cinco
participantes reales, con cobertura de adquisición superior al \pct{97} en
todos los equipos que el modelo emplea y con un horizonte fijo. Son
medidas tomadas en instalaciones en operación, no series sintéticas ni
perfiles tipo.

Lo que definitivamente no se abarcó: el escenario tiene cinco
participantes, y en ningún momento este trabajo se extiende sin más a
comunidades de otro tamaño. El
Capítulo~\ref{sec:esc} muestra que ese número tiene además consecuencias
regulatorias directas, porque algunos umbrales normativos se definen sobre
la participación relativa de cada miembro y con cinco miembros ciertos
umbrales resultan matemáticamente inalcanzables. Tampoco sostiene conclusiones
estacionales completas: nueve meses cubren buena parte del ciclo anual de
Pasto, pero no un año entero.

Hay una segunda ausencia, y es de otra naturaleza: el modelo no representa
la red. No hay impedancias, ni distancias eléctricas, ni límites de
capacidad entre los nodos, de modo que las cinco instituciones entran como
cinco perfiles de carga y de generación sin más relación entre sí que la
temporal.

Esa ausencia no es un descuido sino el alcance del objeto que se simula.
Lo que la regulación colombiana define para una comunidad energética es
una liquidación en kilovatios hora sobre un periodo, no un despacho
físico: el intercambio entre pares se salda a través de la red del
operador, y ninguna de las resoluciones que el Capítulo~\ref{sec:esc}
compara exige un flujo de potencia para calcular lo que cada miembro paga.
El modelo reproduce esa liquidación y no otra cosa.

La consecuencia hay que decirla con el mismo cuidado. Los resultados
describen lo que la liquidación regulatoria arrojaría, no si la red física
podría alojar esos intercambios, y un despliegue real exigiría un estudio
de flujo de carga que este trabajo no hace. De esa familia de restricciones
el documento solo mide la que tiene precio en la factura, que es el consumo
de energía reactiva, y lo hace en el
Anexo~\ref{sec:no-representado}.

En resumen, de las 27 fuentes instrumentadas el modelo emplea 16: los
nueve medidores de frontera, los cinco inversores que definen la generación
y los dos que Udenar aporta a la reconstrucción de su demanda. Sobre esa
base se apoya el tratamiento que describe el capítulo siguiente.
